\documentclass{article}
\usepackage{graphicx}
\usepackage{amsmath}
\usepackage{amssymb}
\usepackage{amsfonts}
\usepackage{hyperref}
\usepackage{exercises}
\usepackage{indentfirst}
\usepackage{amsthm}
\usepackage{csquotes}
\usepackage{yhmath}
\usepackage{comment}
\usepackage{placeins}
\usepackage{tikz}
\newcommand*{\oldneg}{\mathord{\sim}}
\let\oldemptyset\emptyset

\theoremstyle{definition}
\newtheorem{definition}{Definition}

\setlength{\parindent}{.5in}
\setlength{\parskip}{12pt}
\title{Chapter 3 Answers}
\begin{document}
\date{}
\maketitle

The proofs for the propositions 1.1-1.2 actually prove the corresponding propositions for the other three sections. Hence, we only prove the 
theorems in section 1, while the others can be easily seen to follow from these proofs. 

\section*{1.1}
This is a consequence of 1.5 and 1.6, as 1.6 shows that any universal register is creative, while 1.5 shows that any creative register can be 
stumped (and hence any universal register can be stumped).

\section*{1.2}
Suppose that at least one of the registers, call it $R_n$, is universal. Hence, for all numbers $x$ and $y$, $R_n$ affirms $xy$ iff $R_x$ affirms 
$y$. Now, $R_{n^{\#'}}$ affirms $n^{\#'}$ iff $R_n$ affirms $n^{\#'}n^{\#'}$ iff $R_{n^\#}$ affirms $n^{\#'}$ iff $R_{n^{\#'}}$ denies $n^{\#'}$. Thus, 
$R_{n^{\#'}}$ affirms $n^{\#'}$ iff $R_{n^{\#'}}$ defines $n^{\#'}$, thus meaning $R_{n^{\#'}}$ is stumped by $n^{\#'}$, its own index. 

\section*{1.3}
This, like 1.1, is a consequence of 1.5 and 1.6.

\section*{1.4}
(a): Suppose some register $R_n$ is universal. Thus $R_{n^{\#}}$ affirms $n^{\#'}$ iff $R_n$ affirms $n^{\#'}n^{\#'}$ iff $R_{n^{\#'}}$ affirms $n^{\#'}$ 
iff $R_{n^{\#}}$ denies $n^{\#'}$, 
 meaning $R_{n^\#}$ is stumped by $n^{\#'}$, the index of its opposer.


(b): Next, suppose that some register $R_n$ is contra-universal. Hence, for all $x$ and $y$, $R_{n'}$ affirms $xy$ iff $R_n$ denies $xy$ iff 
$R_x$ affirms $y$. Thus $R_{n'}$ is a universal register. Thus, because there is a universal register, we know that some register is stumped by its own 
index by Proposition 1.2. In particular: $R_{n^\#}$ affirms $n^\#$ iff $R_n$ affirms $n^\#n^\#$ iff $R_{n^{\#}}$ denies $n^\#$, 
thus proving that $R_{n^\#}$ is stumped by $n^\#$, its own index. 

\section*{1.5}
Suppose some register $C$ is creative. Thus for every register $R$, there is at least one number $n$ such that $C$ affirms $n$ iff $R$ affirms $n$. Well, 
consider $R_{k'}$ for $k$ the index of $C$. Thus there is at least one number $n$ such that $R_k$ (that is, $C$) affirms $n$ iff 
$R_{k'}$ affirms $n$ iff $R_k$ denies $n$, and hence $R_k$ ($C$) must be stumped by $n$.

\section*{1.6}
(a): First, we show that every universal register is creative. Suppose $U$ is universal, that is, for all numbers $x$ and $y$, $U$ affirms 
$xy$ iff $R_x$ affirms $y$. Now, consider any register $R_n$. We must show that there is at least one number $m$ such that $U$ affirms $m$ iff 
$R_n$ affirms $n$. Well, $U$ affirms $n^\#n^\#$ iff $R_{n^\#}$ affirms $n^\#$ iff $R_n$ affirms $n^\#n^\#$, thus proving that $U$ is creative.


(b): Next, we must show that every contra-universal register is creative. Suppose $V$ is contra-universal, that is, for all numbers $x$ and $y$,
$V$ affirms $xy$ iff $R_x$ denies $y$. Now, consider any register $R_n$. We must show that there is at least one number $m$ such that $U$ affirms $m$ 
iff $R_n$ affirms $m$. Well, $V$ affirms $n^{\#'}n^{\#'}$ iff $R_{n^{\#'}}$ denies $n^{\#'}$ iff $R_{n^\#}$ affirms $n^{\#'}$ iff 
$R_n$ affirms $n^{\#'}n^{\#'}$, thus proving that $V$ is creative. 

\section*{1.7}
Consider some register $R_x$. Now, $R_x$ affirms $x^\#x^\#$ iff $R_{x^\#}$ affirms $x^\#$ (which is to say $x^\#x^\#\in A$) and $R_x$ denies $x^\#x^\#$ iff 
$R_{x^\#}$ denies $x^\#$ (which is to say $x^\#x^\#\in B$). Hence for every register $R_x$, $x^\#x^\#$ is a fixed point of $R_x$.

\section*{1.8}
We suppose that $U$ is universal and $V$ is contra-universal.


(a): Suppose $xy$ is a fixed point of $U'$. Hence, for every $x$, 
$R_x$ affirms $y$ iff $U$ affirms $xy$ iff $U'$ denies $xy$ iff $xy\in B$ iff $R_x$ denies $y$, 
thus meaning $U$ is stumped by $xy$, the fixed point of $U'$.


(b): Suppose $xy$ is a fixed point of $V$. Hence $R_x$ denies $y$ iff $V$ affirms $xy$ iff $xy\in A$ iff $R_x$ affirms $y$, thus meaning $V$ is 
stumped by $xy$, its own fixed point.

\section*{1.9}
Suppose some register $R$ is such that $R$ affirms all number in $A$ and doesn't affirm any number in $B$. Hence, if $xy\in A$, then $R$ affirms 
$xy$, while if $xy\in B$, $A$ doesn't affirm $xy$ (i.e. either denies or is stumped by it). Let $R_k=R$ (that is, let $k$ be the index of $R$). 


(a): On the one hand, let $k^{\#'}k^{\#'}$ be in $A$. Then $R_{k^{\#'}}$ would affirm $k^{\#'}$ and $R_k$ would affirm $k^{\#'}k^{\#'}$, and thus 
$R_{k^\#}$ would deny $k^{\#'}$ and thus $R_k$ would deny $k^{\#'}k^{\#'}$, which is a contradiction. Thus $k^{\#'}k^{\#'}\notin A$.


(b): On the other hand, suppose that $k^{\#'}k^{\#'}\in B$. Then $R_{k^{\#'}}$ would deny $k^{\#'}$ and $R_k$ would not affirm $k^{\#'}$, and thus 
$R_{k^\#}$ would affirm $k^{\#'}$ and thus $R_k$ would affirm $k^{\#'}k^{\#'}$, which is a contradiction. Thus $k^{\#'}k^{\#'}\notin B$.


Hence $R_{k^{\#'}}$, because $k^{\#'}k^{\#'}$ is in neither $A$ nor $B$, must be stumped by $k^{\#'}$. So $R_k$ affirms 
$k^{\#'}k^{\#'}$ iff $R_{k'}$ denies $k^{\#'}k^{\#'}$ iff $R_{k^{\#'}}$ affirms $k^{\#'}k^{\#'}$, meaning $R_k$ can neither affirm nor deny 
$k^{\#'}k^{\#'}$, and hence is stumped by it.


\textbf{Note:} Smullyan said in his answer that we could show that $R_k$ is stumped by $xy$ for $xy$ the fixed point of $R$. However, I was unable to see 
this. 

\section*{1.10}
Suppose that some register $R_k$ is a stump detector. Thence $R_k$ halts for $k^\#k^\#$ iff $R_{k^\#}$ is stumped by $k^\#$ iff 
$R_k$ is stumped by $k^\#k^\#$. Thus $R_k$ halts for $k^\#k^\#$ iff $R_k$ is stumped by $k^\#k^\#$, thus meaning $R_k$ cannot exist.

\section*{1.11}
We suppose that $U$ is universal, $V$ is contra-universal, and that $R$ dominates $U$ and $R'$ dominatnes $V$. 
Now, every number $xy$ in $A$ is such that $R_x$ affirms $y$, meaning $U$ affirms $xy$ and hence $R$ affirms $xy$. Also, any number $xy$ in $B$ is 
such that $R_x$ denies $y$ meaning $V$ affirms $xy$ meaning $R'$ affirms $xy$ and thus $R$ denies $xy$. Hence $R$ affirms all numbers in $A$ and doesn't 
affirm any number in $B$. Thus, by Proposition 1.9, $R$ can be stumped. 

\section*{1.12}
We suppose that there is some universal register $U$. Now, let $S$ be the affirmation set for $U$. Hence, because, for all $x$ and $y$, 
$U$ affirms $xy$ iff $R_x$ affirms $y$, $S$ contains all those and only those numbers $xy$ such that $R_x$ affirms $y$ (hence $S$ is the set $A$). 
Hence $S'$ (that is, the complement of $S$) contains all those and only those numbers $xy$ such that $R_x$ does not affirm $y$. But suppose that there was 
some register $R$ such that $S'$ was $R$'s affirmation set. Now, this would mean that $R$ would affirm all those and only those numbers $xy$ such that 
$R_x$ does not affirm $y$. Let $k$ be the index of $R$. Hence $R_k$ affirms $k^\#k^\#$ iff $R_{k^\#}$ did not affirm $k^\#$ (i.e. either denies or 
is stumped by it) iff $R_k$ did not affirm $k^\#k^\#$, which means such an $R_k$ cannot exist. Therefore $S'$ (that is, $A'$) is not an affirmation set. 


$Alternatively$: We can instead show there cannot be any register $R_k$ (as in the above) using fixed points. Consider some $R_k$ such that $S'$ is 
its affirmation set. Let $xy$ be its fixed point. Thus $R_k$ affirms $xy$ iff $xy\in A$ and $R_k$ denies $xy$ iff $xy\in B$. 
So, if $xy\in A$, $R_k$ affirms $A$, while if $xy\in A'$, $R_k$ does not affirm $xy$. In any case, $S'$ (that is, $A'$) cannot be the affirmation set 
for any register $R_k$, as either $R_k$ affirms some number $xy$ outside of $A'$, or $A'$ contains some number $xy$ such that $R_k$ does not affirm $xy$ 
(both of these cases would mean $S'$ and the set of numbers affirmed by $R_k$ are not the same). 

\end{document}
